% ============================================================
% Theoretical Analysis — FairPoC Fairness Guarantee
% Paste into main.tex as \section{Theoretical Analysis}
% ============================================================

\section{Theoretical Analysis}
\label{sec:theory}

We provide formal guarantees for FairPoC's contribution fairness
and characterise the conditions under which baseline methods fail
to achieve fairness under non-IID data heterogeneity.

% ── Notation ────────────────────────────────────────────────
\subsection{Notation and Problem Formulation}

Let $\mathcal{C} = \{1, \ldots, n\}$ denote the set of participating
clients in a given FL round.
Each client $i$ is assigned a contribution score
$s_i \in \mathbb{R}_{+}$ computed by the PoC formula:
\begin{equation}
  s_i \;=\; \alpha \cdot \mathrm{acc}_i \;+\; \beta \cdot \mathrm{div}_i,
  \label{eq:poc_score}
\end{equation}
where $\mathrm{acc}_i$ is client $i$'s validation accuracy,
$\mathrm{div}_i = \tfrac{1 - \cos(\Delta w_i,\,\bar{\Delta w})}{2}$
is the gradient divergence, and $\alpha + \beta = 1$.
Let $S = \sum_{i=1}^{n} s_i$ and let $R \in \mathbb{R}_{+}$ be the
total token reward pool per round.
Client $i$ receives reward $r_i$ and we define its
\emph{reward-contribution ratio}:
\begin{equation}
  \varphi_i \;=\; \frac{r_i}{s_i}.
\end{equation}
Perfect fairness requires $\varphi_i$ to be equal for all clients:
every client receives reward proportional to its contribution.
We measure this with Jain's fairness index~\cite{jain1984}:
\begin{equation}
  \mathcal{J}(\boldsymbol{\varphi})
  \;=\;
  \frac{\bigl(\sum_{i=1}^{n} \varphi_i\bigr)^2}
       {n \cdot \sum_{i=1}^{n} \varphi_i^2}
  \;\in\; \left(\tfrac{1}{n},\, 1\right].
\end{equation}
$\mathcal{J} = 1$ iff all $\varphi_i$ are equal (perfect proportionality);
$\mathcal{J} = 1/n$ iff one client receives all the reward.

% ── Proposition 1 ───────────────────────────────────────────
\subsection{FairPoC Fairness Guarantee}

\begin{proposition}[Guaranteed Proportional Fairness]
\label{prop:fairness}
Let rewards be allocated by the FairPoC smart contract as
\begin{equation}
  r_i \;=\; \frac{s_i}{S} \cdot R, \quad \forall\, i \in \mathcal{C}.
  \label{eq:proportional}
\end{equation}
Then $\mathcal{J}(\boldsymbol{\varphi}) = 1$ for any scores
$s_i > 0$ and any reward pool $R > 0$.
\end{proposition}

\begin{proof}
Under~\eqref{eq:proportional}:
\[
  \varphi_i
  \;=\; \frac{r_i}{s_i}
  \;=\; \frac{s_i / S \cdot R}{s_i}
  \;=\; \frac{R}{S}, \quad \forall\, i \in \mathcal{C}.
\]
Since $\varphi_i = R/S$ is constant across all clients:
\[
  \mathcal{J}(\boldsymbol{\varphi})
  \;=\; \frac{\bigl(n \cdot \tfrac{R}{S}\bigr)^2}
             {n \cdot n \cdot \bigl(\tfrac{R}{S}\bigr)^2}
  \;=\; \frac{n^2 (R/S)^2}{n^2 (R/S)^2}
  \;=\; 1. \qed
\]
\end{proof}

\begin{remark}[Trustless Enforcement]
Proposition~\ref{prop:fairness} holds regardless of the score
values $s_i$, as long as rewards are distributed
via~\eqref{eq:proportional}.
In FairPoC, equation~\eqref{eq:proportional} is executed by the
\texttt{RewardDistributor} smart contract — an immutable on-chain
program that no participant can modify or bypass.
Therefore, $\mathcal{J} = 1$ is a \emph{trustless guarantee}:
it holds without requiring any client to trust the aggregator
or any other participant.
\end{remark}

% ── Proposition 2 ───────────────────────────────────────────
\subsection{Fairness Degradation of Uniform Baselines}

We next show that uniform reward methods (FedAVG, BLADE-FL)
achieve $\mathcal{J} = 1$ only under the restrictive condition
that all clients have equal contribution — a condition that never
holds under non-IID heterogeneity.

\begin{proposition}[Uniform Reward Fairness Condition]
\label{prop:uniform_fairness}
For any method assigning uniform rewards $r_i = R/n$ for all
$i \in \mathcal{C}$:
\[
  \mathcal{J}(\boldsymbol{\varphi}) = 1
  \;\iff\;
  s_1 = s_2 = \cdots = s_n.
\]
\end{proposition}

\begin{proof}
With $r_i = R/n$, the reward-contribution ratios are
$\varphi_i = R/(n \cdot s_i)$.
By the Cauchy-Schwarz inequality,
$(\sum \varphi_i)^2 \leq n \sum \varphi_i^2$,
with equality iff all $\varphi_i$ are equal, i.e., iff all $s_i$
are equal.
Hence $\mathcal{J}(\boldsymbol{\varphi}) = 1 \iff s_i = s_j$
for all $i, j \in \mathcal{C}$. \qed
\end{proof}

\begin{corollary}[Non-IID Fairness Gap]
\label{cor:noniid_gap}
Under non-IID data heterogeneity (Dirichlet parameter
$\alpha < \infty$), client contribution scores $s_i$ are not
equal in general.
Therefore, uniform reward methods yield $\mathcal{J} < 1$, and
the fairness gap $1 - \mathcal{J}$ grows with the degree of
heterogeneity.
\end{corollary}

Corollary~\ref{cor:noniid_gap} explains the empirical results in
Section~\ref{sec:experiments}: at $\alpha = 0.1$ (extreme non-IID),
FedAVG and BLADE-FL exhibit a statistically significant fairness
gap ($p < 0.01$, Table~\ref{tab:multiseed}), while FairPoC
achieves $\mathcal{J} = 1.000$ with zero variance across all seeds
by Proposition~\ref{prop:fairness}.

% ── Proposition 3 ───────────────────────────────────────────
\subsection{Stability of Gradient Divergence as Incentive Signal}

Shapley-value-based methods — including PoShapley-BCFL~\cite{poshapley2023}
— rely on the marginal accuracy contribution of each client,
which diminishes as the global model converges~\cite{poshapley2023}.
We formalise why gradient divergence avoids this degradation.

\begin{proposition}[Non-Degeneracy of Gradient Divergence]
\label{prop:stability}
Let $\Delta w_i^{(t)} = w_i^{(t)} - w_{\mathrm{global}}^{(t)}$ be
client $i$'s gradient update at round $t$.
For any two clients $i \neq j$ whose local data distributions
$\mathcal{D}_i \neq \mathcal{D}_j$ differ,
there exists $\epsilon > 0$ such that:
\[
  \bigl|\mathrm{div}_i^{(t)} - \mathrm{div}_j^{(t)}\bigr| \;>\; \epsilon
  \quad \forall\, t \geq 1,
\]
provided the local objectives are non-degenerate
(i.e., $\nabla \mathcal{L}_i \neq \nabla \mathcal{L}_j$ at any local
optimum).
\end{proposition}

\begin{proof}[Proof sketch]
Since $\mathcal{D}_i \neq \mathcal{D}_j$, the local loss
functions $\mathcal{L}_i \neq \mathcal{L}_j$.
At any round $t$, local training from $w_{\mathrm{global}}^{(t)}$
produces gradients $\Delta w_i^{(t)}$ aligned with
$-\nabla \mathcal{L}_i(w_{\mathrm{global}}^{(t)})$.
Since $\mathcal{D}_i \neq \mathcal{D}_j$, the gradient directions
$\Delta w_i^{(t)}$ and $\Delta w_j^{(t)}$ differ by at least
$\epsilon > 0$ in cosine distance, independent of the global model's
convergence state.
Hence $\mathrm{div}_i^{(t)}$ remains differentiated from
$\mathrm{div}_j^{(t)}$ throughout training. \qed
\end{proof}

\begin{remark}
Proposition~\ref{prop:stability} confirms that gradient divergence
provides a \emph{stable and non-degenerate incentive signal}
at all training stages — in contrast to marginal accuracy
(LOO or Shapley), which converges to zero as the global model
approaches convergence.
This directly addresses the open problem noted in
PoShapley-BCFL~\cite{poshapley2023}: \emph{``the marginal model
contribution becomes smaller and smaller during the late stage
of model convergence.''}
\end{remark}

% ── Proposition 4 ───────────────────────────────────────────
\subsection{Gradient Norm Degrades as a Fairness Proxy at Extreme Non-IID}

We formalise why gradient norm — used in FairPoC's PoC score —
loses alignment with the data-quality ground truth under
degenerate heterogeneity ($\alpha \to 0$).

\begin{proposition}[Gradient Norm Decouples from Data Volume at Extreme Non-IID]
\label{prop:gradnorm_limit}
Let client $i$ hold $n_i$ samples with class distribution
$\{p_{i,k}\}_{k=1}^{C}$.
Define client $i$'s \emph{class concentration} as
$\rho_i = \max_k p_{i,k}$ (fraction of samples from the dominant class).
Under cross-entropy loss, the gradient update satisfies:
\begin{equation}
  \|\Delta w_i\| \;\approx\;
  \rho_i \cdot \left\|\nabla_w \mathcal{L}\!\left(f(\cdot;\,w),\,k_i^*\right)\right\|,
  \label{eq:gradnorm_approx}
\end{equation}
where $k_i^* = \arg\max_k p_{i,k}$ is the dominant class.
In particular, $\|\Delta w_i\|$ is \emph{independent of $n_i$} up to
the factor $\rho_i$.
\end{proposition}

\begin{proof}
Client $i$ computes the empirical gradient:
\[
  \Delta w_i \;=\; \frac{1}{n_i}
    \sum_{j=1}^{n_i} \nabla_w \mathcal{L}(f(x_j;\,w),\,y_j).
\]
Under extreme non-IID ($\rho_i \to 1$), almost all labels satisfy
$y_j = k_i^*$.
Partitioning the sum by class:
\[
  \Delta w_i \;=\;
    \sum_{k=1}^{C} \hat{p}_{i,k} \cdot
    \underbrace{\frac{1}{n_{i,k}} \sum_{j: y_j=k}
    \nabla_w \mathcal{L}(f(x_j;\,w),\,k)}_{\approx\; g_k(w)},
\]
where $\hat{p}_{i,k} = n_{i,k}/n_i$ is the empirical class fraction
and $g_k(w)$ is the expected gradient for class $k$.
Taking norms and applying the triangle inequality:
\[
  \|\Delta w_i\| \;\leq\;
    \sum_{k} \hat{p}_{i,k} \|g_k(w)\|
  \;\approx\; \rho_i \|g_{k_i^*}(w)\|,
\]
since $\hat{p}_{i,k} \approx 0$ for $k \neq k_i^*$.
The $n_i$ factors cancel in the class-average $g_k(w)$,
so $\|\Delta w_i\|$ depends on $\rho_i$ and the current model
state $w$, but \emph{not} on $n_i$ directly. \qed
\end{proof}

\begin{corollary}[Misalignment with Data-Quality GT at Extreme Non-IID]
\label{cor:gradnorm_misalignment}
The data-quality ground truth is
$\mathrm{gt}_i = \gamma (n_i/N) + (1-\gamma) H_i$,
where $H_i$ is the normalised class entropy.
Under extreme non-IID ($\rho_i \to 1$):
$H_i \to 0$, so $\mathrm{gt}_i \approx \gamma (n_i / N)$,
meaning $\mathrm{gt}_i \propto n_i$.
By Proposition~\ref{prop:gradnorm_limit}, $\|\Delta w_i\|$ is
independent of $n_i$.
Therefore:
\[
  \frac{\|\Delta w_i\|}{\mathrm{gt}_i}
  \;\approx\;
  \frac{\rho_i \|g_{k_i^*}(w)\|}{\gamma\, n_i/N}
  \;\propto\; \frac{1}{n_i},
\]
which varies inversely with client data volume.
Clients with fewer samples ($n_i$ small) receive
disproportionately high rewards relative to their data-quality GT,
causing the Jain fairness index of reward-to-GT ratios to drop
below that of uniform (FedAVG) rewards.
\end{corollary}

\begin{remark}[Scope of the Limitation]
Corollary~\ref{cor:gradnorm_misalignment} applies only under
degenerate extreme non-IID ($\rho_i \approx 1$, i.e.,
Dirichlet $\alpha \ll 1$).
At realistic heterogeneity ($\alpha \geq 0.5$), clients hold
multiple classes ($\rho_i < 1$, $H_i > 0$), and $\|\Delta w_i\|$
correlates with data volume through the multi-class sum in
equation~\eqref{eq:gradnorm_approx}.
This is confirmed empirically: FairPoC achieves the highest
GT-fairness at $\alpha \in \{0.5, 1.0\}$ on both CIFAR-10 and
FashionMNIST (Section~\ref{sec:experiments}).
\end{remark}
